
          %%%%% ~~~~~~~~~~~~~~~~~~~~ %%%%%

\section{Quality indicators}

          %%%%% ~~~~~~~~~~~~~~~~~~~~ %%%%%

% The accuracy of statistical information is the degree to which the information correctly 
% describes the phenomena it was designed to measure.
The quality indicators for this 
product are represented by letters --- ranging from A through E ---
with each letter corresponding to a specific range of values
of coefficient of variation (CV) of the design-based area estimates; see table below.
Most of the inland area of the country is composed of estimates
that are of good, very good or excellent quality.
Estimates with poorer quality (D and E) are generally associated
with rare and uncommon land cover types for a given target geography.

\begin{center}
\begin{tabular}{|c|c|c|}
\hline
$\overset{{\color{white}1}}{\textnormal{Quality indicator}}$
&
Coefficient of variation
&
$\underset{{\color{white}.}}{\textnormal{Description}}$
\\
\hline
$\overset{{\color{white}1}}{\textnormal{A}}$ & $<$ 5.00\% & Excellent 
\\
$\overset{{\color{white}1}}{\textnormal{B}}$ & 5.00\% to 9.99\% & Very good 
\\
$\overset{{\color{white}1}}{\textnormal{C}}$ & 10.00\% to 14.99\% & Good 
\\
$\overset{{\color{white}1}}{\textnormal{D}}$ & 15.00\% to 24.99\% & Acceptable 
\\
$\overset{{\color{white}1}}{\textnormal{E}}$ & $\geq$ 25.00\% & $\underset{{\color{white}.}}{\textnormal{Use with caution}}$
\\
\hline
\end{tabular}
\end{center}

% Table 1 
% Data quality indicators, 2020 land cover area estimates 
% Quality indicator Coefficient of variation value Description 
% A < 5.00% Excellent 
% B 5.00% to 9.99% Very good 
% C 10.00% to 14.99% Good 
% D 15.00% to 24.99% Acceptable 
% E >= 25.00% Use with caution 

          %%%%% ~~~~~~~~~~~~~~~~~~~~ %%%%%
